% BEGIN REV09_MULTIPLICIDAD_MODAL
\subsubsection{Multiplicité et récupération des itinéraires modaux}
\label{corpus9:modal:multiplicidad}

L'incidence \(F_{\mathrm{av}}\), construite par le calendrier tritique
dans le développement précédent, agit ici sur son enveloppe de marches
modales. Nous écrivons
\[
 F=F_{\mathrm{av}}=\begin{pmatrix}0&1\\1&1\end{pmatrix},
 \qquad 0=\Sigma,\quad 1=\Pi.
\]
La coordonnée d'auto-échelle \(\varphi\) et la demi-droite positive
\(r=(1,\varphi)\) sont les sorties déjà obtenues de cette incidence.
Sur celles-ci, on détermine la multiplicité de la lecture scalaire et
on construit une représentation réversible de chaque itinéraire.

On considère l'ensemble des marches dirigées finies
\(\gamma=(x_0,\ldots,x_n)\) de \(F\). Leurs extrémités
et leurs longueurs sont distinguées, les répétitions sont admises et les marches triviales sont incluses.
Ce domaine est l'enveloppe modale ; l'état complet du TPK
conserve en outre les registres de son histoire enrichie.
Le caractère multiplicatif prend la forme
\begin{equation}
 \chi(\gamma)=\varphi^n\frac{r_{x_0}}{r_{x_n}}
             =\varphi^{k(\gamma)},\qquad
 r=(1,\varphi),\quad k(\gamma)=n+x_0-x_n.
 \label{corpus9:modal:caracter}
\end{equation}
Les transitions \(\Sigma\to\Pi\), \(\Pi\to\Pi\) et
\(\Pi\to\Sigma\) incrémentent \(k\) de \(0\), \(1\) et \(2\),
respectivement. La première étape modifie la route tout en maintenant son
caractère. En particulier, \((1,1)\) et \((0,1,1)\) produisent tous deux
\(\varphi\) sous cette lecture.

\begin{proposition}[Multiplicité de la lecture d'auto-échelle]
\label{corpus9:modal:conteo}
Soient \(\mathcal F_k=\{\gamma:k(\gamma)=k\}\) et
\(N_k=|\mathcal F_k|\). Alors
\begin{equation}
 N_0=3,\qquad
 N_k=2\operatorname{tr}(F^k)
     =2\bigl(\varphi^k+(-\varphi^{-1})^k\bigr)
     \quad(k\geq1),
 \label{corpus9:modal:multiplicidad-exacta}
\end{equation}
et
\[
 \lim_{k\to\infty}N_k^{1/k}=\varphi.
\]
\end{proposition}

\begin{proof}
Les extrémités \(i,j\in\{0,1\}\) étant données, la longueur est déterminée par \(n=k-i+j\). Le nombre de marches ayant ces extrémités est \((F^n)_{ij}\). Pour \(k\geq1\), la somme des quatre secteurs donne
\[
 N_k=\operatorname{tr}(F^k)
      +(F^{k+1})_{01}+(F^{k-1})_{10}.
\]
L'identité \(F^2=F+I\) implique, par récurrence, que pour \(n\geq1\)
\[
 F^n=\begin{pmatrix}f_{n-1}&f_n\\f_n&f_{n+1}\end{pmatrix},
 \qquad f_0=0,\quad f_1=1,\quad f_{n+1}=f_n+f_{n-1}.
\]
Pour \(k\geq2\), les deux termes non diagonaux sont
\(f_{k+1}+f_{k-1}=\operatorname{tr}(F^k)\) ; pour \(k=1\),
leur somme est \(1+0=1=\operatorname{tr}(F)\), en utilisant \(F^0=I\).
Pour \(k=0\), les seules marches sont \((0)\), \((1)\) et \((0,1)\).
Les valeurs propres de \(F\) sont \(\varphi\) et \(-\varphi^{-1}\),
de sorte que sa trace fournit la formule fermée.
Enfin,
\[
 N_k=2\varphi^k\bigl(1+(-\varphi^{-2})^k\bigr),
\]
et le facteur entre parenthèses tend vers un ; en prenant les racines
\(k\)-ièmes, on obtient la limite.
\end{proof}

L'auto-échelle coïncide ainsi avec le taux exponentiel de croissance
des généalogies modales finies que le caractère scalaire
rend indiscernables.

\paragraph{Raffinement par composition des retours.}
Soient \(b(\gamma)\) le nombre d'arêtes \(\Pi\to\Sigma\) et \(c(\gamma)\) celui des arêtes \(\Pi\to\Pi\). Les incréments de \eqref{corpus9:modal:caracter} donnent \(k=2b+c\). Avec \(z,t\) comme indéterminées formelles, l'incidence graduée
\[
 M(z,t)=\begin{pmatrix}0&1\\tz^2&z\end{pmatrix}
\]
enregistre simultanément l'exposant et les retours. L'unique arête
de degré zéro est \(\Sigma\to\Pi\), et elle ne possède pas de cycle de degré zéro.
Chaque coefficient compte donc un ensemble fini de marches.
Si \(N_{k,b}\) désigne leur multiplicité de valeurs \(k,b\), la somme
sur toutes les longueurs satisfait
\begin{equation}
 \sum_{k,b\geq0}N_{k,b}z^kt^b
 =\mathbf1^{\mathsf T}(I-M(z,t))^{-1}\mathbf1
 =\frac{3-z+tz^2}{1-z-tz^2},
 \qquad \mathbf1=(1,1)^{\mathsf T}.
 \label{corpus9:modal:serie-graduada}
\end{equation}
En effet, l'entrée \(ij\) de \(M^n\) compte les marches de
longueur \(n\) ayant ces extrémités et leurs poids. La somme formelle
coefficient par coefficient fournit \((I-M)^{-1}\).
Le déterminant de \(I-M\) est \(1-z-tz^2\), et la somme des
quatre entrées de sa comatrice transposée est \(3-z+tz^2\).

Pour \(k\geq1\) et \(0\leq b\leq\lfloor k/2\rfloor\), il vient
\begin{equation}
 N_{k,b}=\frac{2k}{k-b}\binom{k-b}{b}.
 \label{corpus9:modal:conteo-retornos}
\end{equation}
Pour le démontrer, le développement
\[
 (1-z-tz^2)^{-1}=\sum_{m\geq0}(z+tz^2)^m
\]
donne le coefficient \(\binom{k-b}{b}\) dans \(z^kt^b\). La multiplication par le numérateur de \eqref{corpus9:modal:serie-graduada} produit
\[
 3\binom{k-b}{b}-\binom{k-b-1}{b}
   +\binom{k-b-1}{b-1}
 =2\binom{k-b}{b}+2\binom{k-b-1}{b-1}.
\]
Les coefficients hors plage sont pris égaux à zéro. Pour \(b\geq1\), l'égalité \(\binom{k-b-1}{b-1}=\frac{b}{k-b}\binom{k-b}{b}\) donne \eqref{corpus9:modal:conteo-retornos} ; pour \(b=0\), les deux membres valent deux. Le terme constant est séparé : \(N_{0,0}=3\). Pour \(k=4\), il y a quatorze marches, réparties en \(2\), \(8\) et \(4\) selon qu'elles contiennent zéro, un ou deux retours. Le même caractère admet donc des compositions internes distinctes.

\begin{proposition}[Récupération réversible et transport des codes]
\label{corpus9:modal:recuperador}
Chaque marche de l'enveloppe modale admet un code
\[
 \mathcal C(\gamma)=(k,a),\qquad 0\leq a<N_k,
\]
avec un décodage \(\mathcal D\) tel que
\[
 \mathcal D\mathcal C=\mathrm{id},\qquad
 \mathcal C\mathcal D=\mathrm{id}.
\]
Si \(A_v(\gamma)=\gamma v\) prolonge par une arête admissible et
\(T\) supprime le dernier état d'une marche de longueur positive,
leurs représentations
\[
 \mathcal E_v=\mathcal C A_v\mathcal D,\qquad
 \widehat T=\mathcal C T\mathcal D
\]
satisfont, dans leurs domaines respectifs,
\begin{equation}
 \widehat T\mathcal E_v=\mathrm{id},\qquad
 k(A_v\gamma)-k(\gamma)=1+x_n-v.
 \label{corpus9:modal:transporte-codigos}
\end{equation}
\end{proposition}

\begin{proof}
Pour chaque \(k\), on ordonne les secteurs non vides selon \((n,i,j)\)
et, au sein de chaque secteur, les marches dans l'ordre lexicographique de leurs
états. Le nombre de continuations de longueur \(m\) de
\(v\) à \(j\) est \((F^m)_{vj}\). En parcourant une marche, on additionne
les tailles des secteurs précédents et, à chaque position, celles
des continuations lexicographiquement antérieures à celle choisie.
Le résultat est son indice \(a\), compté à partir de zéro.

Pour inverser le processus, on soustrait de l'indice la taille de chaque
secteur précédent jusqu'à localiser l'unique secteur qui le contient.
À chaque position, on procède de même avec les tailles des
continuations. Ces ensembles forment une partition disjointe
et exhaustive du secteur de départ. Pour la longueur zéro, l'unique
marche d'un secteur non vide est son sommet. En supposant récupérées
les continuations de longueur \(m-1\), les intervalles d'indices
de chaque premier successeur identifient de manière univoque ce successeur et l'
indice restant de sa continuation. La récurrence sur \(m\) prouve
que les deux procédures sont inverses, aussi bien sur les marches que
sur les indices \(0,\ldots,N_k-1\).

L'ordre lexicographique sérialise la route existante ; les opérations qui génèrent la route précèdent cette sérialisation. Puisque \(T A_v=\mathrm{id}\), les deux identités inverses donnent
\[
 \widehat T\mathcal E_v
 =\mathcal C T\mathcal D\mathcal C A_v\mathcal D
 =\mathcal C T A_v\mathcal D
 =\mathrm{id}.
\]
Enfin, la définition de l'exposant donne
\[
 k(A_v\gamma)=(n+1)+x_0-v,\qquad
 k(\gamma)=n+x_0-x_n,
\]
dont la différence est la deuxième identité.
\end{proof}

\paragraph{Capacité de distinction avec exposant connu.}
Lorsque \(k\) est connu, une représentation binaire injective de longueur
fixe \(\ell\) exige \(2^\ell\geq N_k\), par le principe des
tiroirs. Par conséquent,
\(\ell\geq\lceil\log_2N_k\rceil\).
L'indice \(a\) précédent atteint cette longueur au moyen de sa
représentation binaire, complétée de zéros initiaux si
nécessaire. Ce minimum correspond à la transmission de l'indice
avec \(k\) connu ; la transmission de \(k\) et la délimitation des
codes de longueur variable ont leurs propres coûts. De plus,
\[
 \lim_{k\to\infty}\frac{\log_2N_k}{k}=\log_2\varphi.
\]
L'égalité mesure la capacité combinatoire de distinction des routes dans l'enveloppe modale.

\paragraph{Information de fibre sous une distribution explicite.}
\(k\) étant fixé, soit \(R_k\) une marche choisie uniformément dans \(\mathcal F_k\). Nous utilisons les logarithmes naturels et \(H(R_k)=-\sum_{\gamma\in\mathcal F_k}p_\gamma\log p_\gamma\), avec \(p_\gamma=1/N_k\). Toutes ces marches ont pour caractère \(\varphi^k\), de sorte que
\[
 H(R_k\mid\chi(R_k))=\log N_k,\qquad
 H(R_k\mid\mathcal C(R_k))=0.
\]
La première égalité provient du fait que la lecture est constante sur la fibre ; la seconde, du décodage unique déjà démontré. En particulier, \(\log N_k/k\to\log\varphi\).

Pour \(k=4\), le caractère laisse une incertitude de \(\log14\)
nats. Les trois classes \(B=b(R_4)\) ont les probabilités
\(2/14\), \(8/14\) et \(4/14\). Conditionnée à chaque classe, la
distribution reste uniforme sur ses \(2\), \(8\) ou \(4\)
éléments, respectivement. Par définition de l'entropie
conditionnelle,
\[
 H(R_4\mid B)
 =\frac2{14}\log2+\frac8{14}\log8+\frac4{14}\log4.
\]
L'indice au sein de la classe, conjointement avec \(B\), retrouve univoquement
la marche. Ces valeurs appartiennent à la distribution uniforme
de la fibre modale déclarée. La mesure de l'espace complet des
histoires, la réalisation temporelle de \(k\) et toute lecture
thermique conservent leurs domaines et transducteurs respectifs.
% END REV09_MULTIPLICIDAD_MODAL
